% Options for packages loaded elsewhere
\PassOptionsToPackage{unicode}{hyperref}
\PassOptionsToPackage{hyphens}{url}
\PassOptionsToPackage{dvipsnames,svgnames,x11names}{xcolor}
\documentclass[
  11pt,
]{article}
\usepackage{xcolor}
\usepackage[margin=0.5in]{geometry}
\usepackage{amsmath,amssymb}
\setcounter{secnumdepth}{5}
\usepackage{iftex}
\ifPDFTeX
  \usepackage[T1]{fontenc}
  \usepackage[utf8]{inputenc}
  \usepackage{textcomp} % provide euro and other symbols
\else % if luatex or xetex
  \usepackage{unicode-math} % this also loads fontspec
  \defaultfontfeatures{Scale=MatchLowercase}
  \defaultfontfeatures[\rmfamily]{Ligatures=TeX,Scale=1}
\fi
\usepackage{lmodern}
\ifPDFTeX\else
  % xetex/luatex font selection
  \setmainfont[]{Helvetica}
  \setmonofont[]{Menlo}
\fi
% Use upquote if available, for straight quotes in verbatim environments
\IfFileExists{upquote.sty}{\usepackage{upquote}}{}
\IfFileExists{microtype.sty}{% use microtype if available
  \usepackage[]{microtype}
  \UseMicrotypeSet[protrusion]{basicmath} % disable protrusion for tt fonts
}{}
\makeatletter
\@ifundefined{KOMAClassName}{% if non-KOMA class
  \IfFileExists{parskip.sty}{%
    \usepackage{parskip}
  }{% else
    \setlength{\parindent}{0pt}
    \setlength{\parskip}{6pt plus 2pt minus 1pt}}
}{% if KOMA class
  \KOMAoptions{parskip=half}}
\makeatother
\usepackage{color}
\usepackage{fancyvrb}
\newcommand{\VerbBar}{|}
\newcommand{\VERB}{\Verb[commandchars=\\\{\}]}
\DefineVerbatimEnvironment{Highlighting}{Verbatim}{commandchars=\\\{\}}
% Add ',fontsize=\small' for more characters per line
\usepackage{framed}
\definecolor{shadecolor}{RGB}{248,248,248}
\newenvironment{Shaded}{\begin{snugshade}}{\end{snugshade}}
\newcommand{\AlertTok}[1]{\textcolor[rgb]{0.94,0.16,0.16}{#1}}
\newcommand{\AnnotationTok}[1]{\textcolor[rgb]{0.56,0.35,0.01}{\textbf{\textit{#1}}}}
\newcommand{\AttributeTok}[1]{\textcolor[rgb]{0.13,0.29,0.53}{#1}}
\newcommand{\BaseNTok}[1]{\textcolor[rgb]{0.00,0.00,0.81}{#1}}
\newcommand{\BuiltInTok}[1]{#1}
\newcommand{\CharTok}[1]{\textcolor[rgb]{0.31,0.60,0.02}{#1}}
\newcommand{\CommentTok}[1]{\textcolor[rgb]{0.56,0.35,0.01}{\textit{#1}}}
\newcommand{\CommentVarTok}[1]{\textcolor[rgb]{0.56,0.35,0.01}{\textbf{\textit{#1}}}}
\newcommand{\ConstantTok}[1]{\textcolor[rgb]{0.56,0.35,0.01}{#1}}
\newcommand{\ControlFlowTok}[1]{\textcolor[rgb]{0.13,0.29,0.53}{\textbf{#1}}}
\newcommand{\DataTypeTok}[1]{\textcolor[rgb]{0.13,0.29,0.53}{#1}}
\newcommand{\DecValTok}[1]{\textcolor[rgb]{0.00,0.00,0.81}{#1}}
\newcommand{\DocumentationTok}[1]{\textcolor[rgb]{0.56,0.35,0.01}{\textbf{\textit{#1}}}}
\newcommand{\ErrorTok}[1]{\textcolor[rgb]{0.64,0.00,0.00}{\textbf{#1}}}
\newcommand{\ExtensionTok}[1]{#1}
\newcommand{\FloatTok}[1]{\textcolor[rgb]{0.00,0.00,0.81}{#1}}
\newcommand{\FunctionTok}[1]{\textcolor[rgb]{0.13,0.29,0.53}{\textbf{#1}}}
\newcommand{\ImportTok}[1]{#1}
\newcommand{\InformationTok}[1]{\textcolor[rgb]{0.56,0.35,0.01}{\textbf{\textit{#1}}}}
\newcommand{\KeywordTok}[1]{\textcolor[rgb]{0.13,0.29,0.53}{\textbf{#1}}}
\newcommand{\NormalTok}[1]{#1}
\newcommand{\OperatorTok}[1]{\textcolor[rgb]{0.81,0.36,0.00}{\textbf{#1}}}
\newcommand{\OtherTok}[1]{\textcolor[rgb]{0.56,0.35,0.01}{#1}}
\newcommand{\PreprocessorTok}[1]{\textcolor[rgb]{0.56,0.35,0.01}{\textit{#1}}}
\newcommand{\RegionMarkerTok}[1]{#1}
\newcommand{\SpecialCharTok}[1]{\textcolor[rgb]{0.81,0.36,0.00}{\textbf{#1}}}
\newcommand{\SpecialStringTok}[1]{\textcolor[rgb]{0.31,0.60,0.02}{#1}}
\newcommand{\StringTok}[1]{\textcolor[rgb]{0.31,0.60,0.02}{#1}}
\newcommand{\VariableTok}[1]{\textcolor[rgb]{0.00,0.00,0.00}{#1}}
\newcommand{\VerbatimStringTok}[1]{\textcolor[rgb]{0.31,0.60,0.02}{#1}}
\newcommand{\WarningTok}[1]{\textcolor[rgb]{0.56,0.35,0.01}{\textbf{\textit{#1}}}}
\usepackage{longtable,booktabs,array}
\newcounter{none} % for unnumbered tables
\usepackage{calc} % for calculating minipage widths
% Correct order of tables after \paragraph or \subparagraph
\usepackage{etoolbox}
\makeatletter
\patchcmd\longtable{\par}{\if@noskipsec\mbox{}\fi\par}{}{}
\makeatother
% Allow footnotes in longtable head/foot
\IfFileExists{footnotehyper.sty}{\usepackage{footnotehyper}}{\usepackage{footnote}}
\makesavenoteenv{longtable}
\usepackage{graphicx}
\makeatletter
\newsavebox\pandoc@box
\newcommand*\pandocbounded[1]{% scales image to fit in text height/width
  \sbox\pandoc@box{#1}%
  \Gscale@div\@tempa{\textheight}{\dimexpr\ht\pandoc@box+\dp\pandoc@box\relax}%
  \Gscale@div\@tempb{\linewidth}{\wd\pandoc@box}%
  \ifdim\@tempb\p@<\@tempa\p@\let\@tempa\@tempb\fi% select the smaller of both
  \ifdim\@tempa\p@<\p@\scalebox{\@tempa}{\usebox\pandoc@box}%
  \else\usebox{\pandoc@box}%
  \fi%
}
% Set default figure placement to htbp
\def\fps@figure{htbp}
\makeatother
\setlength{\emergencystretch}{3em} % prevent overfull lines
\providecommand{\tightlist}{%
  \setlength{\itemsep}{0pt}\setlength{\parskip}{0pt}}
\usepackage{xcolor}
\usepackage{fvextra}
\DefineVerbatimEnvironment{Highlighting}{Verbatim}{breaklines,breakanywhere,commandchars=\\\{\}}
\usepackage{graphicx}
\usepackage{float}
\usepackage{microtype}
\usepackage{setspace}
\usepackage{bookmark}
\IfFileExists{xurl.sty}{\usepackage{xurl}}{} % add URL line breaks if available
\urlstyle{same}
\hypersetup{
  pdftitle={Urban Tree Methods with the USDA Urban Tree Database},
  pdfauthor={Rebecca Bronstein},
  colorlinks=true,
  linkcolor={blue},
  filecolor={Maroon},
  citecolor={Blue},
  urlcolor={blue},
  pdfcreator={LaTeX via pandoc}}

\title{Urban Tree Methods with the USDA Urban Tree Database}
\author{Rebecca Bronstein}
\date{2026-04-01}

\begin{document}
\maketitle

{
\setcounter{tocdepth}{2}
\tableofcontents
}
\subsection{Data Source}\label{data-source}

This report uses the USDA Forest Service Research Data Archive
publication \textbf{Urban tree database} and, specifically, the local
file \texttt{Data/TS3\_Raw\_tree\_data.csv} included in this folder. The
archive describes the publication as a multi-city urban tree dataset
covering more than 14,000 measured street and park trees across 17
cities in the United States. Source page:
\url{https://www.fs.usda.gov/rds/archive/catalog/RDS-2016-0005}.

\subsection{Why Tree Methods Fit This
Dataset}\label{why-tree-methods-fit-this-dataset}

This dataset mixes ecology, urban form, and management questions in one
table. Tree methods work especially well here because they can combine
species categories, regional labels, and structural measurements without
forcing everything into a single linear relationship.

In this notebook I use tree methods to answer four practical questions:

\begin{enumerate}
\def\labelenumi{\arabic{enumi}.}
\tightlist
\item
  Can tree structure and context tell us whether a record comes from a
  street tree or a more park-like setting?
\item
  Which features matter most when predicting tree height?
\item
  Can a rule-based model explain which trees provide high car shade?
\item
  Does bagging improve the street-versus-park-like classification?
\end{enumerate}

\subsection{Data Snapshot}\label{data-snapshot}

\begin{Shaded}
\begin{Highlighting}[]
\FunctionTok{kable}\NormalTok{(overview, }\AttributeTok{caption =} \StringTok{"Summary of the USDA urban tree dataset used in this report."}\NormalTok{)}
\end{Highlighting}
\end{Shaded}

\begin{longtable}[]{@{}lr@{}}
\caption{Summary of the USDA urban tree dataset used in this
report.}\tabularnewline
\toprule\noalign{}
Measure & Value \\
\midrule\noalign{}
\endfirsthead
\toprule\noalign{}
Measure & Value \\
\midrule\noalign{}
\endhead
\bottomrule\noalign{}
\endlastfoot
Total tree records & 14487 \\
Distinct cities & 17 \\
Distinct climate regions & 17 \\
Street trees & 8987 \\
Non-street trees grouped as park-like & 5500 \\
Distinct tree types & 12 \\
\end{longtable}

\subsection{Classification Tree: Street Trees vs.~Park-Like
Trees}\label{classification-tree-street-trees-vs.-park-like-trees}

The original \texttt{Park/Street} field has a few small categories such
as nursery and regional big tree. To make the result easier to explain,
I collapse the outcome into two classes:

\begin{enumerate}
\def\labelenumi{\arabic{enumi}.}
\tightlist
\item
  \texttt{Street}
\item
  \texttt{Park-like}
\end{enumerate}

This lets the tree model act like a site classifier: based on size,
canopy form, and land-use context, what kind of planting environment
does a tree look like it belongs to?

\begin{Shaded}
\begin{Highlighting}[]
\NormalTok{site\_tree }\OtherTok{\textless{}{-}} \FunctionTok{rpart}\NormalTok{(Site.Type }\SpecialCharTok{\textasciitilde{}}\NormalTok{ ., }\AttributeTok{data =}\NormalTok{ train\_site, }\AttributeTok{method =} \StringTok{"class"}\NormalTok{)}
\FunctionTok{plot}\NormalTok{(site\_tree, }\AttributeTok{uniform =} \ConstantTok{TRUE}\NormalTok{, }\AttributeTok{margin =} \FloatTok{0.08}\NormalTok{)}
\FunctionTok{text}\NormalTok{(site\_tree, }\AttributeTok{use.n =} \ConstantTok{TRUE}\NormalTok{, }\AttributeTok{cex =} \FloatTok{0.7}\NormalTok{)}
\end{Highlighting}
\end{Shaded}

\pandocbounded{\includegraphics[keepaspectratio]{usda_urban_tree_methods_files/figure-latex/classification_tree-1.pdf}}

\begin{Shaded}
\begin{Highlighting}[]
\NormalTok{site\_pred }\OtherTok{\textless{}{-}} \FunctionTok{predict}\NormalTok{(site\_tree, test\_site, }\AttributeTok{type =} \StringTok{"class"}\NormalTok{)}
\NormalTok{site\_metrics }\OtherTok{\textless{}{-}} \FunctionTok{classification\_metrics}\NormalTok{(test\_site}\SpecialCharTok{$}\NormalTok{Site.Type, site\_pred)}
\NormalTok{site\_importance }\OtherTok{\textless{}{-}} \FunctionTok{top\_importance}\NormalTok{(site\_tree)}

\FunctionTok{kable}\NormalTok{(site\_metrics, }\AttributeTok{caption =} \StringTok{"Classification performance for predicting site type."}\NormalTok{)}
\end{Highlighting}
\end{Shaded}

\begin{longtable}[]{@{}rr@{}}
\caption{Classification performance for predicting site
type.}\tabularnewline
\toprule\noalign{}
Accuracy & Kappa \\
\midrule\noalign{}
\endfirsthead
\toprule\noalign{}
Accuracy & Kappa \\
\midrule\noalign{}
\endhead
\bottomrule\noalign{}
\endlastfoot
0.987 & 0.971 \\
\end{longtable}

\begin{Shaded}
\begin{Highlighting}[]
\FunctionTok{kable}\NormalTok{(site\_importance, }\AttributeTok{caption =} \StringTok{"Most influential variables in the site{-}type tree."}\NormalTok{)}
\end{Highlighting}
\end{Shaded}

\begin{longtable}[]{@{}lr@{}}
\caption{Most influential variables in the site-type
tree.}\tabularnewline
\toprule\noalign{}
Variable & Importance \\
\midrule\noalign{}
\endfirsthead
\toprule\noalign{}
Variable & Importance \\
\midrule\noalign{}
\endhead
\bottomrule\noalign{}
\endlastfoot
Region & 1171.33 \\
LandUse & 645.09 \\
Leaf & 26.31 \\
WireConf & 24.95 \\
Setback & 15.25 \\
Age & 15.03 \\
\end{longtable}

This is a useful planning interpretation of the dataset. The strongest
variables are region, land use, and setback, which suggests that site
context matters at least as much as raw tree dimensions when
distinguishing street trees from trees in park-like settings.

\subsection{Regression Tree: What Predicts Tree
Height?}\label{regression-tree-what-predicts-tree-height}

Urban tree height is shaped by both biology and management. A regression
tree is useful here because it can split on crown height, diameter, and
leaf area in a way that feels close to how an arborist might reason
through the problem.

\begin{Shaded}
\begin{Highlighting}[]
\NormalTok{height\_tree }\OtherTok{\textless{}{-}} \FunctionTok{rpart}\NormalTok{(TreeHt }\SpecialCharTok{\textasciitilde{}}\NormalTok{ ., }\AttributeTok{data =}\NormalTok{ train\_height, }\AttributeTok{method =} \StringTok{"anova"}\NormalTok{)}
\FunctionTok{plot}\NormalTok{(height\_tree, }\AttributeTok{uniform =} \ConstantTok{TRUE}\NormalTok{, }\AttributeTok{margin =} \FloatTok{0.08}\NormalTok{)}
\FunctionTok{text}\NormalTok{(height\_tree, }\AttributeTok{use.n =} \ConstantTok{TRUE}\NormalTok{, }\AttributeTok{cex =} \FloatTok{0.7}\NormalTok{)}
\end{Highlighting}
\end{Shaded}

\pandocbounded{\includegraphics[keepaspectratio]{usda_urban_tree_methods_files/figure-latex/regression_tree-1.pdf}}

\begin{Shaded}
\begin{Highlighting}[]
\NormalTok{height\_pred }\OtherTok{\textless{}{-}} \FunctionTok{predict}\NormalTok{(height\_tree, test\_height)}
\NormalTok{height\_metrics }\OtherTok{\textless{}{-}} \FunctionTok{regression\_metrics}\NormalTok{(test\_height}\SpecialCharTok{$}\NormalTok{TreeHt, height\_pred)}
\NormalTok{height\_importance }\OtherTok{\textless{}{-}} \FunctionTok{top\_importance}\NormalTok{(height\_tree)}

\FunctionTok{kable}\NormalTok{(height\_metrics, }\AttributeTok{caption =} \StringTok{"Regression performance for predicting tree height in meters."}\NormalTok{)}
\end{Highlighting}
\end{Shaded}

\begin{longtable}[]{@{}rr@{}}
\caption{Regression performance for predicting tree height in
meters.}\tabularnewline
\toprule\noalign{}
RMSE & MAE \\
\midrule\noalign{}
\endfirsthead
\toprule\noalign{}
RMSE & MAE \\
\midrule\noalign{}
\endhead
\bottomrule\noalign{}
\endlastfoot
2.22 & 1.62 \\
\end{longtable}

\begin{Shaded}
\begin{Highlighting}[]
\FunctionTok{kable}\NormalTok{(height\_importance, }\AttributeTok{caption =} \StringTok{"Most influential variables in the tree{-}height model."}\NormalTok{)}
\end{Highlighting}
\end{Shaded}

\begin{longtable}[]{@{}lr@{}}
\caption{Most influential variables in the tree-height
model.}\tabularnewline
\toprule\noalign{}
Variable & Importance \\
\midrule\noalign{}
\endfirsthead
\toprule\noalign{}
Variable & Importance \\
\midrule\noalign{}
\endhead
\bottomrule\noalign{}
\endlastfoot
CrnHt & 124011.20 \\
Leaf & 65971.88 \\
AvgCdia & 56974.27 \\
DBH & 50560.95 \\
Age & 32469.10 \\
TreeType & 15566.46 \\
\end{longtable}

The importance scores usually put crown height, leaf area, average crown
diameter, and DBH near the top. That makes the model easy to justify:
urban tree height is strongly tied to overall crown architecture, not
just one diameter measurement.

\subsection{Rule-Based Tree: Which Trees Deliver High Car
Shade?}\label{rule-based-tree-which-trees-deliver-high-car-shade}

Shade is one of the clearest urban ecosystem services in this dataset.
To make that idea more concrete, I create a binary target called
\texttt{Big.Shade}, where trees with higher \texttt{CarShade} values are
labeled \texttt{High\ shade}.

The rule model asks a planning-style question: which combinations of
region and tree structure tend to produce stronger shade benefits?

\begin{Shaded}
\begin{Highlighting}[]
\NormalTok{shade\_rules }\OtherTok{\textless{}{-}} \FunctionTok{C5.0}\NormalTok{(Big.Shade }\SpecialCharTok{\textasciitilde{}}\NormalTok{ ., }\AttributeTok{data =}\NormalTok{ train\_shade, }\AttributeTok{rules =} \ConstantTok{TRUE}\NormalTok{)}
\NormalTok{shade\_pred }\OtherTok{\textless{}{-}} \FunctionTok{predict}\NormalTok{(shade\_rules, test\_shade)}
\NormalTok{shade\_metrics }\OtherTok{\textless{}{-}} \FunctionTok{classification\_metrics}\NormalTok{(test\_shade}\SpecialCharTok{$}\NormalTok{Big.Shade, shade\_pred)}
\NormalTok{shade\_usage }\OtherTok{\textless{}{-}} \FunctionTok{as.data.frame}\NormalTok{(}\FunctionTok{C5imp}\NormalTok{(shade\_rules, }\AttributeTok{metric =} \StringTok{"usage"}\NormalTok{)) }\SpecialCharTok{\%\textgreater{}\%}
  \FunctionTok{rownames\_to\_column}\NormalTok{(}\AttributeTok{var =} \StringTok{"Variable"}\NormalTok{) }\SpecialCharTok{\%\textgreater{}\%}
  \FunctionTok{rename}\NormalTok{(}\AttributeTok{Usage =}\NormalTok{ Overall) }\SpecialCharTok{\%\textgreater{}\%}
  \FunctionTok{arrange}\NormalTok{(}\FunctionTok{desc}\NormalTok{(Usage)) }\SpecialCharTok{\%\textgreater{}\%}
  \FunctionTok{mutate}\NormalTok{(}\AttributeTok{Usage =} \FunctionTok{round}\NormalTok{(Usage, }\DecValTok{2}\NormalTok{))}

\FunctionTok{kable}\NormalTok{(shade\_metrics, }\AttributeTok{caption =} \StringTok{"Rule{-}based performance for predicting high car shade."}\NormalTok{)}
\end{Highlighting}
\end{Shaded}

\begin{longtable}[]{@{}rr@{}}
\caption{Rule-based performance for predicting high car
shade.}\tabularnewline
\toprule\noalign{}
Accuracy & Kappa \\
\midrule\noalign{}
\endfirsthead
\toprule\noalign{}
Accuracy & Kappa \\
\midrule\noalign{}
\endhead
\bottomrule\noalign{}
\endlastfoot
0.935 & 0.703 \\
\end{longtable}

\begin{Shaded}
\begin{Highlighting}[]
\FunctionTok{kable}\NormalTok{(}\FunctionTok{head}\NormalTok{(shade\_usage, }\DecValTok{6}\NormalTok{), }\AttributeTok{caption =} \StringTok{"Most{-}used predictors in the shade rules."}\NormalTok{)}
\end{Highlighting}
\end{Shaded}

\begin{longtable}[]{@{}lr@{}}
\caption{Most-used predictors in the shade rules.}\tabularnewline
\toprule\noalign{}
Variable & Usage \\
\midrule\noalign{}
\endfirsthead
\toprule\noalign{}
Variable & Usage \\
\midrule\noalign{}
\endhead
\bottomrule\noalign{}
\endlastfoot
Region & 100 \\
Park.Street & 0 \\
TreeType & 0 \\
DBH & 0 \\
TreeHt & 0 \\
AvgCdia & 0 \\
\end{longtable}

\begin{Shaded}
\begin{Highlighting}[]
\FunctionTok{summary}\NormalTok{(shade\_rules)}
\end{Highlighting}
\end{Shaded}

\begin{verbatim}
## 
## Call:
## C5.0.formula(formula = Big.Shade ~ ., data = train_shade, rules = TRUE)
## 
## 
## C5.0 [Release 2.07 GPL Edition]      Wed Apr  1 21:36:56 2026
## -------------------------------
## 
## Class specified by attribute `outcome'
## 
## Read 3501 cases (8 attributes) from undefined.data
## 
## Rules:
## 
## Rule 1: (400/77, lift 6.0)
##  Region in {NMtnPr, Piedmt}
##  ->  class High shade  [0.806]
## 
## Rule 2: (3101/150, lift 1.1)
##  Region in {CenFla, GulfCo, InlEmp, InlVal, InterW, LoMidW, MidWst,
##                    NoCalC, NoEast, PacfNW, SacVal, SoCalC, SWDsrt, TpIntW,
##                    Tropic}
##  ->  class Low shade  [0.951]
## 
## Default class: Low shade
## 
## 
## Evaluation on training data (3501 cases):
## 
##          Rules     
##    ----------------
##      No      Errors
## 
##       2  227( 6.5%)   <<
## 
## 
##     (a)   (b)    <-classified as
##    ----  ----
##     323   150    (a): class High shade
##      77  2951    (b): class Low shade
## 
## 
##  Attribute usage:
## 
##  100.00% Region
## 
## 
## Time: 0.0 secs
\end{verbatim}

This section is intentionally more interpretive than predictive. Even
when the rules are short, they help translate the dataset into a
management idea: some regional and structural profiles are much more
likely to correspond to trees that deliver stronger shade.

\subsection{Bagged Trees: Does the Site Classifier Become More
Stable?}\label{bagged-trees-does-the-site-classifier-become-more-stable}

A single classification tree is easy to explain, but it can change if
the sample shifts. Bagging gives us a way to ask whether the
street-versus-park-like pattern is robust or whether it depends too
heavily on one tree.

\begin{Shaded}
\begin{Highlighting}[]
\NormalTok{bag\_model }\OtherTok{\textless{}{-}} \FunctionTok{bagging}\NormalTok{(Site.Type }\SpecialCharTok{\textasciitilde{}}\NormalTok{ ., }\AttributeTok{data =}\NormalTok{ train\_site, }\AttributeTok{nbagg =} \DecValTok{25}\NormalTok{)}
\NormalTok{bag\_pred }\OtherTok{\textless{}{-}} \FunctionTok{predict}\NormalTok{(bag\_model, test\_site)}

\ControlFlowTok{if}\NormalTok{ (}\FunctionTok{is.list}\NormalTok{(bag\_pred) }\SpecialCharTok{\&\&} \StringTok{"class"} \SpecialCharTok{\%in\%} \FunctionTok{names}\NormalTok{(bag\_pred)) \{}
\NormalTok{  bag\_pred }\OtherTok{\textless{}{-}}\NormalTok{ bag\_pred}\SpecialCharTok{$}\NormalTok{class}
\NormalTok{\}}

\NormalTok{bag\_pred }\OtherTok{\textless{}{-}} \FunctionTok{factor}\NormalTok{(bag\_pred, }\AttributeTok{levels =} \FunctionTok{levels}\NormalTok{(test\_site}\SpecialCharTok{$}\NormalTok{Site.Type))}

\NormalTok{bag\_metrics }\OtherTok{\textless{}{-}} \FunctionTok{classification\_metrics}\NormalTok{(test\_site}\SpecialCharTok{$}\NormalTok{Site.Type, bag\_pred) }\SpecialCharTok{\%\textgreater{}\%}
  \FunctionTok{mutate}\NormalTok{(}\AttributeTok{Model =} \StringTok{"Bagged trees"}\NormalTok{)}

\NormalTok{single\_tree\_metrics }\OtherTok{\textless{}{-}}\NormalTok{ site\_metrics }\SpecialCharTok{\%\textgreater{}\%}
  \FunctionTok{mutate}\NormalTok{(}\AttributeTok{Model =} \StringTok{"Single tree"}\NormalTok{)}

\NormalTok{comparison }\OtherTok{\textless{}{-}} \FunctionTok{bind\_rows}\NormalTok{(single\_tree\_metrics, bag\_metrics) }\SpecialCharTok{\%\textgreater{}\%}
  \FunctionTok{select}\NormalTok{(Model, Accuracy, Kappa)}

\FunctionTok{kable}\NormalTok{(comparison, }\AttributeTok{caption =} \StringTok{"Comparing a single site{-}type tree with a bagged ensemble."}\NormalTok{)}
\end{Highlighting}
\end{Shaded}

\begin{longtable}[]{@{}lrr@{}}
\caption{Comparing a single site-type tree with a bagged
ensemble.}\tabularnewline
\toprule\noalign{}
Model & Accuracy & Kappa \\
\midrule\noalign{}
\endfirsthead
\toprule\noalign{}
Model & Accuracy & Kappa \\
\midrule\noalign{}
\endhead
\bottomrule\noalign{}
\endlastfoot
Single tree & 0.987 & 0.971 \\
Bagged trees & 0.996 & 0.991 \\
\end{longtable}

If the bagged model improves only a little, that is still informative.
It means the site signal in this USDA dataset is already strong enough
that a single tree captures much of the story, while the ensemble mostly
adds stability.

\subsection{Conclusion}\label{conclusion}

This USDA urban tree dataset is especially well suited to tree-based
methods because it combines species, structure, region, and land-use
information in one place.

The models here suggest four useful takeaways:

\begin{enumerate}
\def\labelenumi{\arabic{enumi}.}
\tightlist
\item
  Site context can be inferred very accurately from a tree's regional
  and structural profile.
\item
  Tree height behaves like a crown-architecture problem as much as a
  diameter problem.
\item
  Rule-based models make ecosystem-service questions, like shade, easier
  to explain in plain language.
\item
  Bagging confirms that the street-versus-park-like distinction is a
  stable pattern, not a fragile one.
\end{enumerate}

That makes this archive a strong example of how tree methods can support
both ecological interpretation and urban forestry planning.

\end{document}
